% Topic 1.2.08 — Mean-Square Ergodic Processes.

\section{Эргодичность в среднем квадратичном}
\label{sec:1.2.08-ergodicity}

\begin{ticket}
Эргодические по математическому ожиданию (дисперсии, корреляционной функции) в среднем квадратичном случайные процессы. Критерий и достаточное условие эргодичности по математическому ожиданию в среднем квадратичном.
\end{ticket}

\subsection*{Теория}

Усиленный закон больших чисел (\cref{thm:slln}) утверждает, что временно\'{о}е среднее последовательности независимых одинаково распределённых величин сходится к среднему по ансамблю. Для зависимых процессов аналогичный вопрос звучит так: когда среднее по одной траектории восстанавливает параметры распределения по ансамблю? При втором моменте наиболее простой ответ даёт подход в среднем квадратичном.

\medskip
\noindent\textbf{Стационарные в широком смысле процессы.}

\begin{definition}[Стационарный в широком смысле процесс]
\label{def:wss}
Случайный процесс $X = (X_t)_{t \in T}$ с $T \in \{\Z, \R, [0, \infty)\}$ и $X_t \in L^2$ при всех~$t$ называется \emph{стационарным в широком смысле}, если
\begin{enumerate}[label=(\alph*)]
  \item $\E X_t = m$ не зависит от~$t$;
  \item $\Cov(X_s, X_t)$ зависит только от смещения $\tau = t - s$.
\end{enumerate}
Функцию $K(\tau) = \Cov(X_t, X_{t + \tau})$ называют \emph{(авто)\-ко\-ва\-ри\-а\-ци\-он\-ной функцией}, а в случае центрированного процесса~--- \emph{корреляционной функцией}.
\end{definition}

По~\cref{prop:cov-properties} имеем $K(0) = \Var X_t \geq 0$ и $K(-\tau) = K(\tau)$; кроме того,~$K$ как ядро неотрицательно определена и потому допускает спектральное представление~\cite[гл.~VI, \S~1]{Sevastyanov1982}.

\medskip
\noindent\textbf{Эргодичность в среднем квадратичном.} Для стационарного в широком смысле процесса оценкой математического ожидания по одной траектории служит временн\'{о}е среднее
\[
  \bar X_T \;=\; \frac{1}{T} \int_0^{T} X_t\, dt \quad \text{(непрерывное время)}, \qquad
  \bar X_n \;=\; \frac{1}{n} \sum_{k=1}^{n} X_k \quad \text{(дискретное время)}.
\]
Ниже разбирается случай непрерывного времени; дискретный получается заменой интегралов на суммы.

\begin{definition}[Эргодичность в среднем квадратичном]
\label{def:ms-ergodic}
Стационарный в широком смысле процесс~$X$ называется \emph{эргодическим по математическому ожиданию в среднем квадратичном}, если $\bar X_T \to m$ в среднем квадратичном при $T \to \infty$, то есть $\E (\bar X_T - m)^2 \to 0$. Аналогично~$X$ называется эргодическим \emph{по дисперсии} (соответственно \emph{по корреляционной функции}), если временные оценки $\Var X_t$ (соответственно $K(\tau)$ при фиксированном смещении~$\tau$) сходятся в~$L^2$ к соответствующим параметрам ансамбля.
\end{definition}

\begin{theorem}[Критерий эргодичности по математическому ожиданию]
\label{thm:ms-ergodic-criterion}
Пусть~$X$~--- стационарный в широком смысле процесс с математическим ожиданием~$m$ и ковариационной функцией~$K$. Тогда~$X$ эргодичен по математическому ожиданию в среднем квадратичном тогда и только тогда, когда
\[
  \frac{1}{T^2} \int_0^{T} \int_0^{T} K(t - s)\, ds\, dt \;\to\; 0 \qquad (T \to \infty).
\]
Эквивалентным образом (после замены $\tau = t - s$) условие записывается как
\[
  \frac{2}{T} \int_0^{T} \Bigl(1 - \frac{\tau}{T}\Bigr) K(\tau)\, d\tau \;\to\; 0.
\]
\end{theorem}
\proofof{ms-ergodic-criterion}

\begin{corollary}[Достаточное условие]
\label{cor:ms-ergodic-sufficient}
Если $K(\tau) \to 0$ при $\tau \to \infty$, то~$X$ эргодичен по математическому ожиданию в среднем квадратичном. В частности, это выполнено всякий раз, когда $\int_0^{\infty} |K(\tau)|\, d\tau < \infty$.
\end{corollary}
\proofof{ms-ergodic-sufficient}

\begin{remark}
Критерий выделяет роль дальних корреляций. Если $K(\tau)$ убывает к нулю, то дисперсия временн\'{о}го среднего обращается в нуль со скоростью $O(1/T)$ (или быстрее)~--- траектории ``самоусредняются''. Если же $K(\tau)$ остаётся положительной (например, $X_t = Y$ при единственной случайной величине~$Y$, тогда $K(\tau) \equiv \Var Y$), то временн\'{о}е среднее концентрируется не у~$m$, а у самой случайной величины~$Y$, и эргодичность нарушается. Эргодичность по дисперсии и по корреляционной функции сводится к тому же критерию, применённому к вспомогательным процессам $(X_t - m)^2$ и $(X_t - m)(X_{t + \tau} - m)$, при условии существования соответствующих четвёртых моментов~\cite[гл.~VI, \S~5]{Sevastyanov1982}.
\end{remark}

\subsection*{Доказательства}

\begin{proof}[\Cref{thm:ms-ergodic-criterion}]
\proofanchor{ms-ergodic-criterion}
По теореме Фубини и линейности математического ожидания (\cref{prop:exp-properties}\,(a)) $\E \bar X_T = \frac{1}{T} \int_0^T \E X_t\, dt = m$. Поэтому
\[
  \E (\bar X_T - m)^2 = \Var \bar X_T = \frac{1}{T^2} \int_0^T \!\!\int_0^T \Cov(X_s, X_t)\, ds\, dt = \frac{1}{T^2} \int_0^T \!\!\int_0^T K(t - s)\, ds\, dt,
\]
где в последнем равенстве использована стационарность. Значит, сходимость $\bar X_T \to m$ в~$L^2$ эквивалентна стремлению этого двойного интеграла к нулю, что и есть условие из формулировки.

Эквивалентность с однократным интегралом получается заменой переменных $\tau = t - s$, $\sigma = s$ на квадрате $[0, T] \times [0, T]$. Образ~--- множество $\{(\tau, \sigma) : 0 \leq \sigma \leq T,\ -\sigma \leq \tau \leq T - \sigma\}$; при фиксированном~$\tau$ переменная~$\sigma$ пробегает отрезок $[\max(0, -\tau), \min(T, T - \tau)]$ длины $T - |\tau|$ при $|\tau| \leq T$. С учётом симметрии $K(-\tau) = K(\tau)$ интеграл преобразуется в
\[
  \frac{1}{T^2} \int_{-T}^{T} (T - |\tau|) K(\tau)\, d\tau = \frac{2}{T} \int_0^T \Bigl(1 - \frac{\tau}{T}\Bigr) K(\tau)\, d\tau. \qedhere
\]
\end{proof}

\begin{proof}[\Cref{cor:ms-ergodic-sufficient}]
\proofanchor{ms-ergodic-sufficient}
Пусть $K(\tau) \to 0$ при $\tau \to \infty$. По заданному $\varepsilon > 0$ найдётся~$M$, при котором $|K(\tau)| < \varepsilon$ для $\tau \geq M$. При $T > M$ разобьём интеграл из~\cref{thm:ms-ergodic-criterion}:
\[
  \frac{2}{T} \int_0^T \Bigl(1 - \frac{\tau}{T}\Bigr) K(\tau)\, d\tau = \frac{2}{T} \int_0^M (\cdot)\, d\tau + \frac{2}{T} \int_M^T (\cdot)\, d\tau.
\]
Первый интеграл оценивается сверху как $2M \cdot \max_{[0, M]} |K|/T \to 0$ при $T \to \infty$. Второй ограничен величиной $\frac{2}{T} \int_M^T (1 - \tau/T) \varepsilon\, d\tau \leq \varepsilon$ (вес $1 - \tau/T \leq 1$, длина промежутка не превосходит~$T$). Итак, верхний предел не превосходит~$\varepsilon$ при произвольном $\varepsilon > 0$, то есть интеграл стремится к нулю, и~\cref{thm:ms-ergodic-criterion} даёт эргодичность в среднем квадратичном. Случай абсолютной интегрируемости $\int_0^{\infty} |K| < \infty$ ещё проще: $\frac{2}{T} \int_0^T (1 - \tau/T) |K(\tau)|\, d\tau \leq \frac{2}{T} \int_0^{\infty} |K|\, d\tau \to 0$.
\end{proof}

\subsection*{Примеры}

\begin{example}[Белый шум]
\label{ex:white-noise}
Пусть $(X_n)_{n \in \Z}$~--- независимы и одинаково распределены с математическим ожиданием~$m$ и дисперсией~$\sigma^2$. Тогда $K(0) = \sigma^2$ и $K(\tau) = 0$ при $\tau \neq 0$. Очевидно, $K(\tau) \to 0$, так что по~\cref{cor:ms-ergodic-sufficient} процесс эргодичен по математическому ожиданию в среднем квадратичном. Дискретный аналог дисперсии среднего~--- $\Var \bar X_n = \sigma^2/n \to 0$~--- это в точности~\cref{ex:binomial-moments}, переписанный иначе.
\end{example}

\begin{example}[Процесс Орнштейна--Уленбека]
\label{ex:ou-ergodic}
Гауссовский процесс с непрерывным временем из~\cref{prob:fdd-marginal} с $K(\tau) = e^{-|\tau|}$ удовлетворяет $K(\tau) \to 0$ экспоненциально быстро, поэтому он эргодичен по математическому ожиданию в среднем квадратичном. Количественная форма достаточного условия даёт скорость
\[
  \Var \bar X_T = \frac{2}{T} \int_0^T \bigl(1 - \tau/T\bigr) e^{-\tau}\, d\tau = \frac{2}{T}\bigl(1 - (1 - e^{-T})/T\bigr) \sim \frac{2}{T},
\]
асимптотически ту же скорость~$1/T$, что и для белого шума: длинные корреляции достаточно слабы, чтобы не замедлять самоусреднение более чем на константу.
\end{example}

\begin{example}[Постоянный случайный процесс~--- неэргодический]
\label{ex:constant-process}
Пусть~$Y$~--- произвольная случайная величина с $\Var Y = \sigma^2 > 0$, и $X_t \equiv Y$ для всех~$t$. Тогда $\E X_t = \E Y = m$ и $K(\tau) = \Cov(X_s, X_{s + \tau}) = \Var Y = \sigma^2$ при любом смещении~$\tau$. Критерий не выполнен: $\frac{1}{T^2} \int_0^T \int_0^T K(t - s)\, ds\, dt = \sigma^2 \neq 0$. И действительно: $\bar X_T = Y$ при всех~$T$, временн\'{о}е среднее тождественно совпадает с~$Y$ и концентрируется у случайной величины~$Y$, а не у детерминированного значения~$m$.
\end{example}

\begin{problem}
\label{prob:periodic-ergodic}
Пусть $X_t = A \cos(\omega t + \Phi)$ при детерминированных $A, \omega \in \R$ и $\Phi \sim \mathrm{Uniform}(0, 2\pi)$. Покажите, что~$X$ стационарен в широком смысле, найдите его корреляционную функцию и определите, эргодичен ли~$X$ по математическому ожиданию в среднем квадратичном.
\end{problem}

\begin{proof}[Решение]
$\E X_t = A\, \E \cos(\omega t + \Phi) = \frac{A}{2\pi} \int_0^{2\pi} \cos(\omega t + \varphi)\, d\varphi = 0$, поскольку косинус интегрируется по полному периоду к нулю. Для ковариации
\[
  \E (X_s X_t) = A^2 \E \cos(\omega s + \Phi)\cos(\omega t + \Phi) = \tfrac{A^2}{2}\, \E\bigl(\cos(\omega(t - s)) + \cos(\omega(s + t) + 2\Phi)\bigr).
\]
Второе слагаемое после взятия математического ожидания обнуляется (фаза $2\Phi$ равномерно распределена на $(0, 4\pi)$), и остаётся $\Cov(X_s, X_t) = K(t - s) = \frac{A^2}{2} \cos(\omega(t - s))$. Оба условия из~\cref{def:wss} выполнены (математическое ожидание~--- константа~$0$, ковариация зависит только от смещения), значит,~$X$ стационарен в широком смысле.

Что касается эргодичности, $K(\tau) = \frac{A^2}{2} \cos(\omega \tau)$ \emph{не} стремится к нулю, и достаточное условие непригодно. Применим~\cref{thm:ms-ergodic-criterion} непосредственно: при $\omega \neq 0$
\[
  \frac{2}{T} \int_0^T \bigl(1 - \tau/T\bigr) \tfrac{A^2}{2} \cos(\omega \tau)\, d\tau = \frac{A^2}{T} \cdot \frac{1 - \cos(\omega T)}{\omega^2 T} = O(1/T^2) \to 0,
\]
так что процесс \emph{эргодичен} по математическому ожиданию в среднем квадратичном, несмотря на то что~$K$ не убывает: осцилляции усредняются к нулю. (При $\omega = 0$ мы возвращаемся к~\cref{ex:constant-process}, и эргодичность нарушается.) Вывод: критерий из~\cref{thm:ms-ergodic-criterion} тоньше достаточного условия из~\cref{cor:ms-ergodic-sufficient}.
\end{proof}
